
\subsubsection{6.1 Summary of Findings}

The two models produced different error type distributions among their failed generations. The CodeRL model showed:
\begin{itemize}
\item A non-zero proportion of assertion errors (2.12\% vs. 0\%)
\item Greater mean execution depth before failure (29.4\% vs. 0.09\%)
\item More distributed errors across fine-grained categories
\end{itemize}

These differences were statistically significant with effect sizes ranging from small-medium (V = 0.21 for coarse error types) to large (d = 1.69 for execution depth; V = 0.82 for fine-grained taxonomy).

\subsubsection{6.2 Interpretation Limitations}

The observed differences cannot be attributed solely to alignment method differences. The models differ in:

\begin{enumerate}
\item \textbf{Architecture:} CodeRL uses an encoder-decoder architecture (CodeT5), while CodeLlama-Instruct uses a decoder-only architecture (Llama 2). Architectural differences may affect code generation patterns independent of training objectives.
\end{enumerate}

\begin{enumerate}
\item \textbf{Model size:} CodeRL has 770M parameters; CodeLlama-Instruct has 7B parameters. Despite having fewer parameters, CodeRL achieved much higher pass rates (56.5\% vs. 2.2\%), suggesting that model size is not the primary factor in these results.
\end{enumerate}

\begin{enumerate}
\item \textbf{Pre-training data:} CodeT5 was pre-trained primarily on code, while Llama 2 was pre-trained primarily on general text. This may affect code generation quality independent of fine-tuning.
\end{enumerate}

\begin{enumerate}
\item \textbf{Fine-tuning procedure:} CodeRL received execution-based RL training specifically for code generation. CodeLlama-Instruct received general instruction tuning, not code-specialized DPO. A fairer comparison would use a model explicitly trained with code-focused DPO.
\end{enumerate}

\begin{enumerate}
\item \textbf{Base model performance:} The large gap in pass rates (56.5\% vs. 2.2\%) suggests the models differ substantially in overall code generation capability, which may confound error type comparisons.
\end{enumerate}

\subsubsection{6.3 What These Results Do and Do Not Show}

\textbf{These results show:} Two specific models---CodeRL and CodeLlama-7B-Instruct---produce different error type distributions on HumanEval+ and MBPP+. The differences are statistically significant and the effect sizes are non-trivial.

\textbf{These results do not show:} That execution-based RL training causes different error distributions than preference-based alignment. The confounds enumerated above prevent causal attribution.

\subsubsection{6.4 The Zero-Reward Basin Hypothesis}

The original hypothesis proposed that binary execution reward creates a ''zero-reward basin'' where all non-executing code receives identical zero reward, creating pressure to achieve syntactic validity first. The observed pattern---where CodeRL failures occur deeper in execution and include more assertion errors---is consistent with this hypothesis but does not confirm it, given the confounds.

An alternative explanation is that CodeT5's encoder-decoder architecture and code-focused pre-training simply produce more syntactically valid code than CodeLlama's decoder-only architecture with general-text pre-training, independent of RL training.

\subsubsection{6.5 Effect Amplification at Fine-Grained Taxonomy}

The increase in effect size from V = 0.21 to V = 0.82 at fine-grained taxonomy occurred because CodeLlama-Instruct errors concentrated almost exclusively in syntax\_error (99.8\%), while CodeRL errors distributed primarily between indentation\_error (69.5\%) and syntax\_error (22.9\%). This concentration vs. distribution pattern drives the larger chi-square statistic at the fine-grained level.
